\section{Truncation and construction over the prime field}
\label{sec:construction}

The order bound gives a finite amount of information in the variable $z$ that
detects every nonzero formula.  We first express the operators on this finite
space as matrices over $\F_p(t)$.  We then retain bounded Taylor expansions in
$t$ and represent multiplication by these truncated polynomials as matrices
over $\F_p$.  A common denominator bound ensures that the second truncation
preserves the nonvanishing established by the first.

\subsection{Truncating the word operators}

Fix a prime $p$ and integers $n,s\geq 1$, and set
\begin{equation}\label{eq:z-truncation-length}
  N=1+n\binom{2s}{2}.
\end{equation}
Let $F$ be any field of characteristic $p$ and put $K=F(t)$, with $t$ an
indeterminate over $F$.  Each operator $\mathcal J_i$ raises $z$-order by at
least one, so it induces an operator on $K[[z]]/(z^N)$.  In the ordered basis
$1,z,\ldots,z^{N-1}$, with vectors represented by coefficient columns, its
matrix is
\begin{equation}\label{eq:truncated-operators}
 (U_i)_{r,q}=
 \begin{cases}
 \displaystyle\frac{(1+t)^{i(r-q-1)}}{t^r-1},&0\leq q<r<N,\\[6pt]
 0,&0\leq r\leq q<N,
 \end{cases}
 \qquad 1\leq i\leq n.
\end{equation}
Indeed, the coefficient of $z^r$ in
$b_i(z)z^q=z^{q+1}/(1-(1+t)^i z)$ is $(1+t)^{i(r-q-1)}$ for $r>q$, and
inverting $\sigma-1$ divides that coefficient by $t^r-1$.  In particular,
$U_i\in\Mat_N(\F_p(t))$, independently of the choice of $F$.

Let $e_0$ be the coefficient column of the constant series $1$.  The recursion
$h_{iw}=\mathcal J_i(h_w)$ gives
\begin{equation}\label{eq:operator-evaluation}
 f(U_1,\ldots,U_n)e_0
   =\text{the coefficient column of }h_f\bmod z^N
\end{equation}
for every $f\in F\langle x_1,\ldots,x_n\rangle$.  The order of the operators
agrees with the order of matrix multiplication: the rightmost operator of a
word acts first.

Suppose that $f$ is nonzero and has a formula of size at most $s$.  Word-series
independence gives $h_f\ne0$.  Section~\ref{sec:systems} represents this series
as $uHv$ in a system of dimension $m\leq2s$ with numerator and denominator
degrees at most $n$.  Theorem~\ref{thm:order} therefore gives
\[
  \ord_z h_f\leq n\binom m2\leq n\binom{2s}{2}=N-1.
\]
Thus \eqref{eq:operator-evaluation} is nonzero, and in particular
\begin{equation}\label{eq:rational-hitting}
 f(U_1,\ldots,U_n)\ne0.
\end{equation}
The tuple $U$ already has the required simultaneous detection property, but its
entries are rational functions.  We next replace these functions by finite
matrices while preserving this property.

\subsection{Retaining finitely many coefficients of \texorpdfstring{$t$}{t}}

All denominators in \eqref{eq:truncated-operators} are nonzero at $t=0$.
Moreover, one denominator clears every polynomial evaluated at $U$.  Define
\begin{equation}\label{eq:t-truncation-length}
 D(t)=\prod_{r=1}^{N-1}(t^r-1),
 \qquad
 E=1+\binom N2+n(N-1).
\end{equation}
For every $f\in F\langle x_1,\ldots,x_n\rangle$, each entry of $Df(U)$ is a
polynomial in $F[t]$ of degree less than $E$.  To verify this, consider a
nonzero summand in the $(r,q)$ entry of a word product
$U_{i_1}\cdots U_{i_k}$.  Since the matrices are strictly lower triangular,
the summand follows a chain
\[
 q=j_0<j_1<\cdots<j_k=r<N.
\]
Its denominator is $\prod_{a=1}^k(t^{j_a}-1)$, which divides $D$ because the
indices $j_a$ are distinct.  Its numerator is a power of $1+t$ whose exponent
is
\[
 \sum_{a=1}^k i_{k-a+1}(j_a-j_{a-1}-1)
 \leq n(r-q-k)\leq n(N-1).
\]
Multiplying by $D$ therefore gives degree at most
$\deg D+n(N-1)=E-1$.  For the empty word the assertion follows from
$\deg D=\binom N2$; words of length at least $N$ give zero.  Taking
$F$-linear combinations proves the assertion for $f$.  This reasoning does
not require the factors of $D$ to be relatively prime.

We use the following elementary descent statement to turn this degree bound
into matrices over the coefficient field.  For a field $k$, write
\[
 k[t]_{(t)}=
 \left\{\frac{a(t)}{b(t)}:a,b\in k[t],\ b(0)\ne0\right\}.
\]
Reduction modulo $t^E$ defines a $k$-algebra map from this ring to
$k[t]/(t^E)$, since a polynomial with nonzero constant term is a unit in the
quotient.

\begin{lemma}[Descent by a finite coefficient ring]\label{lem:coefficient-descent}
 Let $k$ be a field, let $n,N,E\geq1$, and let
 $V_1,\ldots,V_n\in\Mat_N(k[t]_{(t)})$.  Reduce their entries modulo $t^E$,
 and replace each resulting entry by its matrix of multiplication on
 $k[t]/(t^E)$ in the basis $1,t,\ldots,t^{E-1}$.  This produces matrices
 $T_1,\ldots,T_n\in\Mat_{NE}(k)$.

 Let $F/k$ be any field extension and
 $f\in F\langle x_1,\ldots,x_n\rangle$.  Suppose that a polynomial
 $D\in F[t]$ with $D(0)\ne0$ satisfies
 \[
  Df(V)\in\Mat_N(F[t]),
  \qquad \deg_t\bigl((Df(V))_{a,b}\bigr)<E
  \quad\text{whenever }(Df(V))_{a,b}\ne0.
 \]
 If $f(V)\ne0$ over $F(t)$, then $f(T)\ne0$ over $F$.
\end{lemma}

\begin{proof}
 Put $R_F=F[t]/(t^E)$ and let
 $\pi_F:F[t]_{(t)}\longrightarrow R_F$ be reduction.  Since $f(V)\ne0$,
 some entry of $Df(V)$ is a nonzero polynomial of degree less than $E$.
 Its reduction is nonzero.  Hence
 \[
   \pi_F(D)\,f(\pi_F(V))=\pi_F(Df(V))\ne0,
 \]
 and therefore $f(\pi_F(V))\ne0$.

 Multiplication on $R_F$ gives a unital $F$-algebra map
 \[
  \rho_F:R_F\longrightarrow\Mat_E(F).
 \]
 It is injective: if multiplication by $a\in R_F$ is zero, applying it to
 $1$ gives $a=0$.  Replacing entries by their multiplication matrices thus
 gives an injective unital $F$-algebra map
 \[
  \rho_F^{(N)}:\Mat_N(R_F)\longrightarrow\Mat_{NE}(F).
 \]
 The maps $\pi_F$ and $\rho_F$ extend the corresponding maps over $k$.
 Consequently $\rho_F^{(N)}(\pi_F(V_i))$ is the matrix $T_i$ viewed over $F$,
 and
 \[
   f(T)=\rho_F^{(N)}\bigl(f(\pi_F(V))\bigr)\ne0.
 \]
 Unital $F$-linearity also shows that every scalar $a\in F$ acts as
 $aI_{NE}$ throughout this evaluation.
\end{proof}

Apply the construction in Lemma~\ref{lem:coefficient-descent} with
$k=\F_p$ and $V_i=U_i$, using the parameters in
\eqref{eq:z-truncation-length} and \eqref{eq:t-truncation-length}.  We denote
the resulting tuple by
\begin{equation}\label{eq:prime-field-tuple}
 T_{p,n,s}=(T_1,\ldots,T_n),\qquad d=NE.
\end{equation}
Only multiplication and reduction of functions regular at $t=0$ enter this
step.  The shift automorphism is used before truncation; no shift on
$\F_p[t]/(t^E)$ is needed.

\subsection{The uniform algorithm and the main theorem}

\begin{proof}[Proof of Theorem~\ref{thm:main}]
 We first verify the simultaneous guarantee.  Fix any field $F$ of
 characteristic $p$ and any nonzero polynomial $f$ over $F$ computed by a
 formula of size at most $s$.
 Equation~\eqref{eq:rational-hitting} gives $f(U)\ne0$ over $F(t)$.
 The degree bound following \eqref{eq:t-truncation-length} verifies the
 hypotheses of Lemma~\ref{lem:coefficient-descent}, which gives $f(T)\ne0$.
 The tuple in \eqref{eq:prime-field-tuple} depends only on $p,n,s$.  Thus this
 argument applies to all choices of $F$ and all their scalar coefficients
 with the same prime-field output.

 Here is a finite deterministic algorithm for producing that output.
 Compute the integers $N,E,d$.  Represent elements of $\F_p[t]/(t^E)$ by
 length-$E$ coefficient arrays, keeping each coefficient as its residue in
 $\{0,\ldots,p-1\}$.  For $0\leq q<r<N$, the required denominator inverse
 is explicitly
 \begin{equation}\label{eq:inverse-geometric-series}
  (t^r-1)^{-1}
    =-\sum_{j=0}^{\lfloor(E-1)/r\rfloor}t^{rj}
    \pmod {t^E}.
 \end{equation}
 Indeed, its product with $t^r-1$ is $1-t^{r(\lfloor(E-1)/r\rfloor+1)}$,
 which is $1$ in the quotient.  Form $(1+t)^{i(r-q-1)}$ by repeated
 multiplication by $1+t$, reducing modulo $t^E$ and modulo $p$ after every
 operation, and multiply it by \eqref{eq:inverse-geometric-series}.
 This computes the reduction of every entry in
 \eqref{eq:truncated-operators}; the remaining entries are zero.

 If such an entry is $c_0+c_1t+\cdots+c_{E-1}t^{E-1}$, its multiplication
 block has entry $c_{a-b}$ in row $a$, column $b$ when
 $0\leq b\leq a<E$, and zero when $a<b$.  Thus each block is obtained by
 shifts and truncations of one coefficient array.  Output $d$ and the
 matrices $T_i$ densely, with these blocks in the positions prescribed by
 $U_i$ and with every entry written as its ordinary binary residue.

 We record the cost in the Turing model.  Put
 $L=\lceil\log_2p\rceil$ and $B=n+s+L+1$.  The parameter formulas give
 \[
  N=O(ns^2),\qquad E=O(n^2s^4),\qquad d=O(n^3s^6),
 \]
 and in particular $N=O(B^3)$, $E=O(B^6)$, and $d=O(B^9)$.
 All exponents used above are at most $n(N-1)$.  With schoolbook
 polynomial arithmetic, at most
 \[
  O\bigl((1+nN)E^2\bigr)
 \]
 residue operations suffice per entry of the $N$-by-$N$ matrices,
 including the construction of its multiplication block.  For all
 $nN^2$ entries this is $O(B^{23})$ residue operations.  Elementary binary
 arithmetic and division with remainder implement each residue operation
 in a number of bit operations polynomial in $L+1$.

 The dense output has bit length
 \[
  O\bigl(\log(d+1)+nd^2(L+1)\bigr)=O(B^{20}),
 \]
 including entry separators and dimensions.  The algorithm stores only
 polynomially many bits: coefficient arrays, matrices, and loop counters
 all have polynomial size in $B$.  On an ordinary Turing machine, even
 locating every array entry by a scan of this storage adds only a
 polynomial factor.  Integer parameter calculations and input processing
 also have polynomial bit cost.  These explicit finite loops therefore
 give absolute constants $C,k$ bounding both running time and output
 length by $CB^k$, as claimed.
\end{proof}

The same construction applies to the acyclic scalar/variable-labelled path
programs of Section~\ref{sec:systems}.  Their significance here is that their
state bound is the number of graph vertices, even when different paths share
subpaths.

\begin{corollary}\label{cor:path-programs}
 There is a deterministic algorithm which, on input
 $(\operatorname{bin}(p),1^n,1^m)$ with $p$ prime and $n,m\geq1$, outputs
 a tuple in $\Mat_d(\F_p)^n$ detecting every nonzero polynomial computed
 by an acyclic scalar/variable-labelled path program with at most $m$
 vertices, simultaneously over every field of characteristic $p$.
 One may take
 \[
  N=1+n\binom m2,\qquad
  E=1+\binom N2+n(N-1),\qquad d=NE=O(n^3m^6).
 \]
 The running time and dense output length are polynomial in
 $n+m+\lceil\log_2p\rceil$.
\end{corollary}

\begin{proof}
 Section~\ref{sec:systems} gives the same strictly upper triangular
 coefficient representation using at most $m$ states.  Thus
 Theorem~\ref{thm:order} yields $\ord_z h_f\leq n\binom m2$ for every
 nonzero path polynomial.  With this value of $N$, define $U_i$ by
 \eqref{eq:truncated-operators} and apply
 Lemma~\ref{lem:coefficient-descent} with $k=\F_p$ and $E$ as displayed
 above.  The same degree bound and bit analysis apply.  When $m=1$,
 every path polynomial is constant and $N=E=d=1$; the output consists of
 zero $1$-by-$1$ matrices, which detect every nonzero constant.
\end{proof}
